% Reconstructed from the compiled lecture-note PDF.
\chapter{Balanced disks and six synchronized boundary curves}
\section{From one boundary point to six}

\begin{displaytext}
Let K = Kmin be a normalized global minimizer. Choose a counterclockwise regular parametrization \
of its boundary, \
t 7−→σ0(t) ∈∂K. \
For each t, the point σ0(t) is the midpoint of a unique critical hexagon. Let the other five edge \
midpoints be \
σ1(t), . . . , σ5(t).
\end{displaytext}

At every time, the six points form a multipoint:


\begin{displaytext}
σj+3(t) = −σj(t), (3.1) \
σj(t) + σj+2(t) + σj+4(t) = 0, (3.2) \
√ \
3 \
det(σj(t), σj+1(t)) = . (3.3) \
2 \
The strict monotonicity of critical hexagons implies that every σj runs counterclockwise around the \
same boundary curve.
\end{displaytext}

\begin{statementbox}{Definition}
Definition 3.1 (Multicurve). A family of curves
\end{statementbox}

\begin{displaytext}
σj : I →R2, j ∈Z/6Z,
\end{displaytext}

\begin{displaytext}
is a multicurve if (σ0(t), . . . , σ5(t)) is a multipoint for every t ∈I.
\end{displaytext}

A multicurve should be visualized as six synchronized particles on the same convex boundary. The parameter t is not six unrelated time variables; it records one moving critical hexagon.


\section{Regularity inherited by all six curves}

\begin{displaytext}
The boundary of a general convex disk can have corners and flat pieces. Reinhardt’s minimizer has \
no corners, so its tangent direction is continuous. If σ0 is a regular C1 parametrization, then all \
other σj are regular C1 curves.
\end{displaytext}

\begin{displaytext}
Here is the mechanism for σ2. Differentiate the constant determinant relation \
√ \
3 \
det(σ0(t), σ2(t)) = . \
2 \
Writing the unit tangent of σ2 as u2(t) and its speed as v2(t) > 0, so that σ′2 = v2u2, gives \
det(σ′0, σ2) + v2 det(σ0, u2) = 0.
\end{displaytext}

The monotonicity triangles guarantee that neither determinant vanishes. Hence


\begin{displaytext}
0(t), σ2(t)) v2(t) = −det(σ′ det(σ0(t), u2(t)),
\end{displaytext}

which is positive and continuous. The same argument cycles through all indices.


\begin{displaytext}
A further geometric estimate bounds the curvature from above by the curvature of the local \
hyperbolic envelopes. If σ0 is parametrized by arclength, its tangent is Lipschitz. The determinant \
relations then transfer Lipschitz regularity to all σ′j. \
Corollary 3.2. Each σ′j is differentiable almost everywhere.
\end{displaytext}

\begin{prooftext}
Proof. A Lipschitz function on an interval is differentiable almost everywhere by Rademacher’s theorem.
\end{prooftext}

This almost-everywhere second derivative is sufficient for curvature controls and optimal-control equations.


\section{The global star condition}

\begin{displaytext}
Fix a time t and abbreviate sj = σj(t). The tangent σ′j(t) cannot point arbitrarily. Convexity \
confines it to the open cone based at sj whose boundary rays point toward
\end{displaytext}

sj+1 and sj + sj+1.


After translating the apex to the origin, this says that σ′j is a positive linear combination of the two edge directions sj+1 −sj, sj+1.


The source calls this the star condition.


Why are the inequalities strict? If a tangent lay exactly on a bounding ray, convexity and the multipoint relations would force another curve to have zero velocity or would create a corner. Both are impossible for a regular minimizing multicurve.


\section{The local curvature condition}

For a regular plane curve γ(t), the signed curvature is


\begin{displaytext}
det(γ′(t), γ′′(t)) \
κsigned(t) = . \
|γ′(t)|3
\end{displaytext}

ray through sj + sj+1


σ′j(t) ray through sj+1


sj


\begin{figureplaceholder}
Figure 3.1: The star condition: the tangent lies strictly inside a convex cone determined by the current critical multipoint.
\end{figureplaceholder}

A counterclockwise simple closed curve bounds a convex disk if and only if its signed curvature is nonnegative, interpreted almost everywhere when the derivative is Lipschitz. Hence


\begin{displaytext}
det(σ′j(t), σ′′j (t)) ≥0
\end{displaytext}

for almost every t.


There are thus two distinct convexity constraints:


• the global star condition, which controls tangent directions relative to the synchronized hexagon; • the local curvature condition, which forbids negative turning.


Both are needed. Local nonnegative curvature alone would not ensure the six curves remain synchronized in the critical-hexagon geometry. The star condition alone would not prevent smallscale nonconvex bending.


\section{Balanced disks}

The paper packages the preceding properties into a class Kbal.


\begin{statementbox}{Definition}
Definition 3.3 (Balanced disk). A centrally symmetric convex disk K is balanced if its boundary admits six regular C1 curves σj such that:
\end{statementbox}

1. the derivatives σ′j are Lipschitz; 2. at each time the six points form a normalized multipoint; 3. the signed curvature is nonnegative almost everywhere; 4. the star condition holds at every time; 5. the six curves have the same image and fit together periodically.


Every normalized global minimizer is balanced. The balanced class may contain non-minimizers; it is a convenient admissible class for the variational problem.


\section{Circle representation}

By affine invariance, choose one critical multipoint and send it to the sixth roots of unity s∗j. Reparametrize time so that this multipoint occurs at t = 0:


\begin{displaytext}
σj(0) = s∗j.
\end{displaytext}

This normalization is called the circle representation. The term does not mean that K is a circle. It means only that one synchronized hexagon is the regular one inscribed in the unit circle.


The balanced Reinhardt problem is now:


minimize area(K) over balanced disks in circle representation.


\section{A first preview of the matrix reduction}

The multipoint equations are linear and preserve determinants. This strongly suggests that every current multipoint is obtained from the standard one by a determinant-one linear map. Indeed, the next chapter proves that there is a unique matrix g(t) ∈SL2(R) such that


\begin{displaytext}
σj(t) = g(t)s∗j (j = 0, . . . , 5).
\end{displaytext}

Thus six plane curves are generated by one curve of matrices. The tangent directions are encoded by a traceless matrix X(t) = g(t)−1g′(t).


The star condition becomes three elementary linear inequalities on the entries of X. The curvature condition becomes a control lying in a triangle.


\begin{insightbox}{Chapter summary}
\end{insightbox}

A minimizing boundary is not treated as one free curve. Every boundary point determines a unique critical hexagon, so six points move together. Their multipoint identities, star inequalities, and nonnegative curvature define a balanced disk. Affine normalization fixes one multipoint as the sixth roots of unity. This makes a matrix representation unavoidable and prepares the passage from convex geometry to a finite-dimensional control system.

